\Article{23}
Bei Versuchen zur Bestimmung des absoluten Werts $T$ des Erdmagnetismus ist es wichtig, alle Versuche innerhalb einer nicht zu langen Zeit abzuschließen, damit keine merkliche Änderung des magnetischen Zustands der verwendeten Nadeln zu befürchten ist. Für die Beobachtung der Ablenkungen der beweglichen Nadel empfiehlt es sich daher, nur die erste Methode aus Artikel $20$ anzuwenden und lediglich \SourceLeaf{116}zwei geeignet gewählte, verschiedene Abstände zu verwenden, sofern zwei Terme der Reihe genügen. Unter mehreren Anwendungen dieser Methode wählen wir hier diejenige als Beispiel, auf die wir die größte Sorgfalt verwendet haben. Die Abstände wurden dabei mit mikroskopischer Genauigkeit gemessen.

Die Versuche wurden am $1832$ Sept. $18$ mit zwei Apparaten durchgeführt, die wir durch die Buchstaben $A,B$ unterscheiden, und mit drei Nadeln, die wir mit den Zahlen $1,2,3$ bezeichnen. Die Nadeln $1,2$ sind dieselben, die in Artikel $11$ als erste und zweite Nadel bezeichnet wurden. Die Versuche gliedern sich in zwei Teile.

Zunächst wurden die gleichzeitigen Schwingungen der Nadel $1$ im Apparat $A$ und der Nadel $2$ im Apparat $B$ beobachtet. Für die Dauer einer Schwingung ergab sich, auf unendlich kleine Amplitude zurückgeführt
\begin{center}
\begin{tabular}{lcl}
Nadel $1$&.\ .\ .\ .\ .\ .&$15''22450$\\
Nadel $2$&.\ .\ .\ .\ .\ .&$17{,}29995$
\end{tabular}
\end{center}
Der erste Wert wurde aus $305$, der zweite aus $264$ Schwingungen ermittelt.

Anschließend wurde Nadel $3$ im Apparat $A$ aufgehängt und Nadel $1$ auf der Geraden senkrecht zum magnetischen Meridian sowohl nach Osten als auch nach Westen, jeweils auf zwei Arten, aufgestellt. Die Ablenkung der Nadel $3$ wurde für jede Lage der Nadel $1$ beobachtet. Die Versuche wurden bei zwei verschiedenen Abständen $R$ wiederholt und ergaben die folgenden Werte des Winkels $v$, der wie in den Artikeln $19$, $21$ zu verstehen ist.
\[
\begin{aligned}
R&=1^{\mathrm m}2,&\qquad v&=3^0\ 42'\ 19''4\\
R'&=1{,}6.&v'&=1\ 34\ 19{,}3
\end{aligned}
\]
Auch während dieser Versuche wurden die Schwingungen der Nadel $2$ im Apparat $B$ beobachtet. Für die mittlere Zeit ergab sich aus $414$ Schwingungen die Dauer einer unendlich kleinen Schwingung $=17''29484$.

Die Zeiten wurden mit einem Chronometer gemessen, dessen täglicher Nachgang $=14''24$ betrug.

Bezeichnen $M,m$ die Momente des freien Magnetismus der Nadeln $1$ und $3$ und $\theta$ die Torsionskonstante des Fadens im Apparat $A$, während er die Nadel $1$ oder $3$ trug (sie haben fast dasselbe Gewicht), so haben wir
\[
\frac{\theta}{TM}=\frac1{597{,}4}
\]
wie in Artikel $11$, und
\[
\frac{\theta}{Tm}=\frac1{721{,}6}
\]
da die Nadel $3$ stärker magnetisiert war als die Nadel $1$.
\par
\SourceLeaf{117}
Das Trägheitsmoment der Nadel $1$ war bereits aus früheren Versuchen bekannt (siehe Artikel $11$). Sie hatten $K=4228732400$ ergeben, mit Millimeter und Milligramm als Einheiten.

Die Thermometeranzeige änderte sich in den beiden Räumen, in denen die Apparate aufgestellt waren, während der gesamten Versuchszeit so wenig, dass man sie nicht berücksichtigen muss.

Gehen wir nun zur Auswertung dieser Versuche über, um daraus die Intensität $T$ des Erdmagnetismus zu bestimmen. Die ungleichen Schwingungszeiten der Nadel $2$ zeigen eine leichte Änderung dieser Intensität. Um einen bestimmten Wert angeben zu können, führen wir daher die beobachtete Schwingungsdauer der Nadel $1$ auf den mittleren Zustand des Erdmagnetismus während des zweiten Teils der Beobachtungen zurück. Diese Zeit erfordert eine weitere Korrektur wegen des Nachgangs des Chronometers und eine dritte wegen der Torsion des Fadens. So ergibt sich für die reduzierte Dauer einer Schwingung der Nadel $1$
\[
\begin{aligned}
&=15{,}22450\times\frac{17{,}29484}{17{,}29995}\mathpunct{.}\frac{86400}{86385{,}76}\mathpunct{.}\BareRoot{\frac{598{,}4}{597{,}4}}\\
&=15''23530=t
\end{aligned}
\]

Daraus ergibt sich der Wert des Produkts $TM=\frac{\pi\pi K}{tt}=179770600$. Der kleine Unterschied zwischen diesem Wert und demjenigen, den wir oben in Artikel $11$ für den $11$. September gefunden haben, ist sowohl der Änderung des Erdmagnetismus als auch der Änderung des magnetischen Zustands der Nadel zuzuschreiben.

Aus den beobachteten Ablenkungen erhalten wir
\[
F=\frac{R'^5\operatorname{tang}v'-R^5\operatorname{tang}v}{R'R'-RR}=113056200
\]
wenn wir den Millimeter als Einheit nehmen. Daraus folgt
\[
\frac MT=\frac12F\left(1+\frac{\theta}{Tm}\right)=56606437
\]
Der Vergleich dieser Zahl mit dem Wert von $TM$ ergibt schließlich
\[
T=1{,}782088
\]
als Wert der Intensität der horizontalen erdmagnetischen Kraft am $18$. September um $5$ Uhr.

